\appendix
\section*{Приложение А. Листинги исходного кода}
\addcontentsline{toc}{section}{Приложение А. Листинги исходного кода}

% ============================================================
% Инфраструктура и конфигурация
% ============================================================

\subsection*{А.1. docker-compose.yml}
\begin{lstlisting}[caption={docker-compose.yml --- PostgreSQL для разработки}, language={}]
services:
  postgres:
    image: postgres:16-alpine
    container_name: idempotency-postgres
    environment:
      POSTGRES_DB: idempotency_demo
      POSTGRES_USER: demo
      POSTGRES_PASSWORD: demo
    ports:
      - "5432:5432"
    volumes:
      - pgdata:/var/lib/postgresql/data
    healthcheck:
      test: ["CMD-SHELL", "pg_isready -U demo -d idempotency_demo"]
      interval: 5s
      timeout: 3s
      retries: 5

volumes:
  pgdata:
\end{lstlisting}

\subsection*{А.2. application.yml}
\begin{lstlisting}[caption={application.yml --- основная конфигурация}, language={}]
spring:
  application:
    name: idempotency-demo
  datasource:
    url: jdbc:postgresql://localhost:5432/idempotency_demo
    username: demo
    password: demo
    driver-class-name: org.postgresql.Driver
  jpa:
    hibernate:
      ddl-auto: validate
    open-in-view: false
    properties:
      hibernate:
        format_sql: true
  flyway:
    enabled: true
    locations: classpath:db/migration
  data:
    rest:
      base-path: /api

app:
  idempotency:
    enabled: true
    key-ttl-hours: 24

server:
  port: 8080
\end{lstlisting}

\subsection*{А.3. application-non-idempotent.yml}
\begin{lstlisting}[caption={application-non-idempotent.yml --- профиль без идемпотентности}, language={}]
app:
  idempotency:
    enabled: false

spring:
  flyway:
    locations: classpath:db/migration,classpath:db/migration-non-idempotent
  jpa:
    properties:
      hibernate:
        '[jakarta.persistence.lock.timeout]': 0
\end{lstlisting}

% ============================================================
% Миграции Flyway
% ============================================================

\subsection*{А.4. V1\_\_init\_schema.sql}
\begin{lstlisting}[caption={V1\_\_init\_schema.sql --- основная схема БД}, language=SQL]
-- Таблица заказов (домен OrderController)
CREATE TABLE orders (
    id UUID PRIMARY KEY DEFAULT gen_random_uuid(),
    description VARCHAR(255) NOT NULL,
    amount DECIMAL(19, 2) NOT NULL,
    status VARCHAR(50) NOT NULL DEFAULT 'CREATED',
    created_at TIMESTAMP NOT NULL DEFAULT NOW(),
    CONSTRAINT uk_order_description_amount UNIQUE (description, amount)
);

-- Таблица платежей (домен PaymentRouter)
CREATE TABLE payments (
    id UUID PRIMARY KEY DEFAULT gen_random_uuid(),
    order_id UUID NOT NULL REFERENCES orders(id),
    amount DECIMAL(19, 2) NOT NULL,
    status VARCHAR(50) NOT NULL DEFAULT 'PENDING',
    created_at TIMESTAMP NOT NULL DEFAULT NOW()
);

-- Таблица товаров (домен ProductRepository)
CREATE TABLE products (
    id BIGSERIAL PRIMARY KEY,
    sku VARCHAR(100) NOT NULL UNIQUE,
    name VARCHAR(255) NOT NULL,
    price DECIMAL(19, 2) NOT NULL,
    quantity INTEGER NOT NULL DEFAULT 0,
    version INTEGER NOT NULL DEFAULT 0
);
\end{lstlisting}

\subsection*{А.5. V2\_\_idempotency\_keys.sql}
\begin{lstlisting}[caption={V2\_\_idempotency\_keys.sql --- таблица ключей идемпотентности}, language=SQL]
CREATE TABLE idempotency_keys (
    id BIGSERIAL PRIMARY KEY,
    idempotency_key VARCHAR(255) NOT NULL UNIQUE,
    method VARCHAR(10) NOT NULL,
    path VARCHAR(500) NOT NULL,
    response_status INTEGER,
    response_body TEXT,
    created_at TIMESTAMP NOT NULL DEFAULT NOW(),
    expires_at TIMESTAMP NOT NULL DEFAULT NOW() + INTERVAL '24 hours'
);

CREATE INDEX idx_idempotency_key ON idempotency_keys(idempotency_key);
CREATE INDEX idx_idempotency_expires ON idempotency_keys(expires_at);
\end{lstlisting}

\subsection*{А.6. V3\_\_drop\_constraints.sql}
\begin{lstlisting}[caption={V3\_\_drop\_constraints.sql --- удаление ограничений (профиль non-idempotent)}, language=SQL]
ALTER TABLE orders DROP CONSTRAINT IF EXISTS uk_order_description_amount;
ALTER TABLE products DROP CONSTRAINT IF EXISTS products_sku_key;
\end{lstlisting}

% ============================================================
% Механизм идемпотентности
% ============================================================

\subsection*{А.7. IdempotencyApplication.java}
\begin{lstlisting}[caption={IdempotencyApplication.java --- точка входа}]
package com.example.idempotency;

import org.springframework.boot.SpringApplication;
import org.springframework.boot.autoconfigure.SpringBootApplication;

@SpringBootApplication
public class IdempotencyApplication {
    public static void main(String[] args) {
        SpringApplication.run(IdempotencyApplication.class, args);
    }
}
\end{lstlisting}

\subsection*{А.8. IdempotencyProperties.java}
\begin{lstlisting}[caption={IdempotencyProperties.java --- настройки идемпотентности}]
package com.example.idempotency.config;

import org.springframework.boot.context.properties.ConfigurationProperties;

@ConfigurationProperties(prefix = "app.idempotency")
public class IdempotencyProperties {
    private boolean enabled = true;
    private int keyTtlHours = 24;

    public boolean isEnabled() { return enabled; }
    public void setEnabled(boolean enabled) { this.enabled = enabled; }
    public int getKeyTtlHours() { return keyTtlHours; }
    public void setKeyTtlHours(int keyTtlHours) { this.keyTtlHours = keyTtlHours; }
}
\end{lstlisting}

\subsection*{А.9. WebMvcConfig.java}
\begin{lstlisting}[caption={WebMvcConfig.java --- регистрация фильтра}]
package com.example.idempotency.config;

import com.example.idempotency.idempotency.IdempotencyFilter;
import com.example.idempotency.idempotency.IdempotencyService;
import org.springframework.boot.context.properties.EnableConfigurationProperties;
import org.springframework.boot.web.servlet.FilterRegistrationBean;
import org.springframework.context.annotation.Bean;
import org.springframework.context.annotation.Configuration;

@Configuration
@EnableConfigurationProperties(IdempotencyProperties.class)
public class WebMvcConfig {
    @Bean
    public FilterRegistrationBean<IdempotencyFilter> idempotencyFilterRegistration(
            IdempotencyService idempotencyService, IdempotencyProperties properties) {
        var reg = new FilterRegistrationBean<IdempotencyFilter>();
        reg.setFilter(new IdempotencyFilter(idempotencyService, properties));
        reg.addUrlPatterns("/api/orders/*", "/api/orders",
                           "/api/payments/*", "/api/payments");
        reg.setOrder(1);
        return reg;
    }
}
\end{lstlisting}

\subsection*{А.10. IdempotencyKeyEntity.java}
\begin{lstlisting}[caption={IdempotencyKeyEntity.java --- сущность ключа идемпотентности}]
package com.example.idempotency.idempotency;

import jakarta.persistence.*;
import java.time.LocalDateTime;

@Entity
@Table(name = "idempotency_keys")
public class IdempotencyKeyEntity {
    @Id @GeneratedValue(strategy = GenerationType.IDENTITY)
    private Long id;

    @Column(name = "idempotency_key", nullable = false, unique = true)
    private String idempotencyKey;

    @Column(nullable = false, length = 10)
    private String method;

    @Column(nullable = false, length = 500)
    private String path;

    @Column(name = "response_status")
    private Integer responseStatus;

    @Column(name = "response_body", columnDefinition = "TEXT")
    private String responseBody;

    @Column(name = "created_at", nullable = false, updatable = false)
    private LocalDateTime createdAt = LocalDateTime.now();

    @Column(name = "expires_at", nullable = false)
    private LocalDateTime expiresAt = LocalDateTime.now().plusHours(24);

    public IdempotencyKeyEntity() {}
    public IdempotencyKeyEntity(String key, String method, String path) {
        this.idempotencyKey = key; this.method = method; this.path = path;
    }

    // Getters и Setters
    public Long getId() { return id; }
    public String getIdempotencyKey() { return idempotencyKey; }
    public Integer getResponseStatus() { return responseStatus; }
    public void setResponseStatus(Integer s) { this.responseStatus = s; }
    public String getResponseBody() { return responseBody; }
    public void setResponseBody(String b) { this.responseBody = b; }
}
\end{lstlisting}

\subsection*{А.11. IdempotencyKeyRepository.java}
\begin{lstlisting}[caption={IdempotencyKeyRepository.java}]
package com.example.idempotency.idempotency;

import org.springframework.data.jpa.repository.JpaRepository;
import java.util.Optional;

public interface IdempotencyKeyRepository
        extends JpaRepository<IdempotencyKeyEntity, Long> {
    Optional<IdempotencyKeyEntity> findByIdempotencyKey(String key);
}
\end{lstlisting}

\subsection*{А.12. IdempotencyService.java}
\begin{lstlisting}[caption={IdempotencyService.java --- сервис управления ключами}]
package com.example.idempotency.idempotency;

import org.springframework.stereotype.Service;
import org.springframework.transaction.annotation.Transactional;
import java.util.Optional;

@Service
public class IdempotencyService {
    private final IdempotencyKeyRepository repository;

    public IdempotencyService(IdempotencyKeyRepository repository) {
        this.repository = repository;
    }

    public Optional<IdempotencyKeyEntity> findByKey(String key) {
        return repository.findByIdempotencyKey(key);
    }

    @Transactional
    public IdempotencyKeyEntity registerKey(String key, String method, String path) {
        return repository.save(new IdempotencyKeyEntity(key, method, path));
    }

    @Transactional
    public void saveResponse(String key, int status, String responseBody) {
        repository.findByIdempotencyKey(key).ifPresent(entity -> {
            entity.setResponseStatus(status);
            entity.setResponseBody(responseBody);
            repository.save(entity);
        });
    }
}
\end{lstlisting}

\subsection*{А.13. IdempotencyFilter.java}
\begin{lstlisting}[caption={IdempotencyFilter.java --- HTTP-фильтр идемпотентности}]
package com.example.idempotency.idempotency;

import com.example.idempotency.config.IdempotencyProperties;
import jakarta.servlet.*;
import jakarta.servlet.http.*;
import org.springframework.dao.DataIntegrityViolationException;
import org.springframework.http.HttpStatus;
import org.springframework.web.filter.OncePerRequestFilter;
import org.springframework.web.util.ContentCachingResponseWrapper;
import java.io.IOException;
import java.nio.charset.StandardCharsets;

public class IdempotencyFilter extends OncePerRequestFilter {
    public static final String IDEMPOTENCY_KEY_HEADER = "Idempotency-Key";
    private final IdempotencyService idempotencyService;
    private final IdempotencyProperties properties;

    public IdempotencyFilter(IdempotencyService svc, IdempotencyProperties props) {
        this.idempotencyService = svc; this.properties = props;
    }

    @Override
    protected void doFilterInternal(HttpServletRequest request,
            HttpServletResponse response, FilterChain chain)
            throws ServletException, IOException {
        if (!"POST".equalsIgnoreCase(request.getMethod())) {
            chain.doFilter(request, response); return;
        }
        if (!properties.isEnabled()) {
            chain.doFilter(request, response); return;
        }
        String key = request.getHeader(IDEMPOTENCY_KEY_HEADER);
        if (key == null || key.isBlank()) {
            response.setStatus(HttpStatus.BAD_REQUEST.value());
            response.setContentType("application/json;charset=UTF-8");
            response.getWriter().write(
                "{\"error\":\"Idempotency-Key required\"}");
            return;
        }
        var existing = idempotencyService.findByKey(key);
        if (existing.isPresent()) {
            var cached = existing.get();
            if (cached.getResponseStatus() != null) {
                response.setStatus(cached.getResponseStatus());
                response.setContentType("application/json;charset=UTF-8");
                if (cached.getResponseBody() != null)
                    response.getWriter().write(cached.getResponseBody());
                return;
            }
            response.setStatus(HttpStatus.CONFLICT.value());
            return;
        }
        try {
            idempotencyService.registerKey(key,
                request.getMethod(), request.getRequestURI());
        } catch (DataIntegrityViolationException e) {
            response.setStatus(HttpStatus.CONFLICT.value()); return;
        }
        var wrapper = new ContentCachingResponseWrapper(response);
        chain.doFilter(request, wrapper);
        String body = new String(
            wrapper.getContentAsByteArray(), StandardCharsets.UTF_8);
        idempotencyService.saveResponse(key, wrapper.getStatus(), body);
        wrapper.copyBodyToResponse();
    }
}
\end{lstlisting}

% ============================================================
% Домен: Заказы (@RestController)
% ============================================================

\subsection*{А.14. Order.java}
\begin{lstlisting}[caption={Order.java --- сущность заказа}]
package com.example.idempotency.order;

import jakarta.persistence.*;
import java.math.BigDecimal;
import java.time.LocalDateTime;
import java.util.UUID;

@Entity @Table(name = "orders")
public class Order {
    @Id @GeneratedValue(strategy = GenerationType.UUID)
    private UUID id;
    @Column(nullable = false)
    private String description;
    @Column(nullable = false, precision = 19, scale = 2)
    private BigDecimal amount;
    @Column(nullable = false) @Enumerated(EnumType.STRING)
    private OrderStatus status = OrderStatus.CREATED;
    @Column(name = "created_at", nullable = false, updatable = false)
    private LocalDateTime createdAt = LocalDateTime.now();

    public Order() {}
    public Order(String description, BigDecimal amount) {
        this.description = description; this.amount = amount;
    }
    // Getters и Setters
    public UUID getId() { return id; }
    public String getDescription() { return description; }
    public void setDescription(String d) { this.description = d; }
    public BigDecimal getAmount() { return amount; }
    public void setAmount(BigDecimal a) { this.amount = a; }
    public OrderStatus getStatus() { return status; }
    public void setStatus(OrderStatus s) { this.status = s; }
}
\end{lstlisting}

\subsection*{А.15. OrderStatus.java}
\begin{lstlisting}[caption={OrderStatus.java}]
package com.example.idempotency.order;

public enum OrderStatus { CREATED, CONFIRMED, CANCELLED }
\end{lstlisting}

\subsection*{А.16. OrderDto.java}
\begin{lstlisting}[caption={OrderDto.java --- DTO заказа с валидацией}]
package com.example.idempotency.order;

import jakarta.validation.constraints.*;
import java.math.BigDecimal;

public record OrderDto(
    @NotBlank(message = "Описание обязательно") String description,
    @NotNull @DecimalMin("0.01") BigDecimal amount,
    OrderStatus status
) {}
\end{lstlisting}

\subsection*{А.17. OrderRepository.java}
\begin{lstlisting}[caption={OrderRepository.java}]
package com.example.idempotency.order;

import org.springframework.data.jpa.repository.JpaRepository;
import java.util.UUID;

public interface OrderRepository extends JpaRepository<Order, UUID> {}
\end{lstlisting}

\subsection*{А.18. OrderService.java}
\begin{lstlisting}[caption={OrderService.java --- бизнес-логика заказов}]
package com.example.idempotency.order;

import org.springframework.stereotype.Service;
import org.springframework.transaction.annotation.Transactional;
import java.util.*;

@Service
public class OrderService {
    private final OrderRepository orderRepository;

    public OrderService(OrderRepository repo) { this.orderRepository = repo; }

    public List<Order> findAll() { return orderRepository.findAll(); }

    public Order findById(UUID id) {
        return orderRepository.findById(id)
            .orElseThrow(() -> new OrderNotFoundException(id));
    }

    @Transactional
    public Order create(OrderDto dto) {
        return orderRepository.save(new Order(dto.description(), dto.amount()));
    }

    @Transactional
    public Order update(UUID id, OrderDto dto) {
        var order = findById(id);
        order.setDescription(dto.description());
        order.setAmount(dto.amount());
        if (dto.status() != null) order.setStatus(dto.status());
        return orderRepository.save(order);
    }

    @Transactional
    public void delete(UUID id) { orderRepository.delete(findById(id)); }

    public static class OrderNotFoundException extends RuntimeException {
        public OrderNotFoundException(UUID id) {
            super("Order not found: " + id);
        }
    }
}
\end{lstlisting}

\subsection*{А.19. OrderController.java}
\begin{lstlisting}[caption={OrderController.java --- тип контроллера \textnumero 1: @RestController}]
package com.example.idempotency.order;

import jakarta.validation.Valid;
import org.springframework.http.*;
import org.springframework.web.bind.annotation.*;
import java.util.*;

@RestController
@RequestMapping("/api/orders")
public class OrderController {
    private final OrderService orderService;

    public OrderController(OrderService svc) { this.orderService = svc; }

    @GetMapping
    public List<Order> list() { return orderService.findAll(); }

    @GetMapping("/{id}")
    public Order get(@PathVariable UUID id) { return orderService.findById(id); }

    @PostMapping
    public ResponseEntity<Order> create(@Valid @RequestBody OrderDto dto) {
        return ResponseEntity.status(HttpStatus.CREATED)
            .body(orderService.create(dto));
    }

    @PutMapping("/{id}")
    public Order update(@PathVariable UUID id, @Valid @RequestBody OrderDto dto) {
        return orderService.update(id, dto);
    }

    @DeleteMapping("/{id}") @ResponseStatus(HttpStatus.NO_CONTENT)
    public void delete(@PathVariable UUID id) { orderService.delete(id); }
}
\end{lstlisting}

% ============================================================
% Домен: Платежи (Functional Endpoints)
% ============================================================

\subsection*{А.20. Payment.java}
\begin{lstlisting}[caption={Payment.java --- сущность платежа}]
package com.example.idempotency.payment;

import com.example.idempotency.order.Order;
import jakarta.persistence.*;
import java.math.BigDecimal;
import java.time.LocalDateTime;
import java.util.UUID;

@Entity @Table(name = "payments")
public class Payment {
    @Id @GeneratedValue(strategy = GenerationType.UUID)
    private UUID id;
    @ManyToOne(fetch = FetchType.EAGER)
    @JoinColumn(name = "order_id", nullable = false)
    private Order order;
    @Column(nullable = false, precision = 19, scale = 2)
    private BigDecimal amount;
    @Column(nullable = false) @Enumerated(EnumType.STRING)
    private PaymentStatus status = PaymentStatus.PENDING;
    @Column(name = "created_at", nullable = false, updatable = false)
    private LocalDateTime createdAt = LocalDateTime.now();

    public Payment() {}
    public Payment(Order order, BigDecimal amount) {
        this.order = order; this.amount = amount;
    }
    // Getters и Setters
    public UUID getId() { return id; }
    public BigDecimal getAmount() { return amount; }
    public PaymentStatus getStatus() { return status; }
}
\end{lstlisting}

\subsection*{А.21. PaymentStatus.java}
\begin{lstlisting}[caption={PaymentStatus.java}]
package com.example.idempotency.payment;

public enum PaymentStatus { PENDING, COMPLETED, FAILED }
\end{lstlisting}

\subsection*{А.22. PaymentDto.java}
\begin{lstlisting}[caption={PaymentDto.java}]
package com.example.idempotency.payment;

import java.math.BigDecimal;
import java.util.UUID;

public record PaymentDto(UUID orderId, BigDecimal amount) {}
\end{lstlisting}

\subsection*{А.23. PaymentRepository.java}
\begin{lstlisting}[caption={PaymentRepository.java}]
package com.example.idempotency.payment;

import org.springframework.data.jpa.repository.JpaRepository;
import java.util.UUID;

public interface PaymentRepository extends JpaRepository<Payment, UUID> {}
\end{lstlisting}

\subsection*{А.24. PaymentService.java}
\begin{lstlisting}[caption={PaymentService.java --- бизнес-логика платежей}]
package com.example.idempotency.payment;

import com.example.idempotency.order.*;
import org.springframework.stereotype.Service;
import org.springframework.transaction.annotation.Transactional;
import java.util.*;

@Service
public class PaymentService {
    private final PaymentRepository paymentRepository;
    private final OrderRepository orderRepository;

    public PaymentService(PaymentRepository pr, OrderRepository or) {
        this.paymentRepository = pr; this.orderRepository = or;
    }

    public List<Payment> findAll() { return paymentRepository.findAll(); }

    public Payment findById(UUID id) {
        return paymentRepository.findById(id)
            .orElseThrow(() -> new PaymentNotFoundException(id));
    }

    @Transactional
    public Payment create(PaymentDto dto) {
        var order = orderRepository.findById(dto.orderId())
            .orElseThrow(() -> new OrderService.OrderNotFoundException(dto.orderId()));
        return paymentRepository.save(new Payment(order, dto.amount()));
    }

    public static class PaymentNotFoundException extends RuntimeException {
        public PaymentNotFoundException(UUID id) {
            super("Payment not found: " + id);
        }
    }
}
\end{lstlisting}

\subsection*{А.25. PaymentHandler.java}
\begin{lstlisting}[caption={PaymentHandler.java --- тип контроллера \textnumero 2: HandlerFunction}]
package com.example.idempotency.payment;

import com.fasterxml.jackson.databind.ObjectMapper;
import org.springframework.stereotype.Component;
import org.springframework.web.servlet.function.*;
import java.net.URI;
import java.util.UUID;

@Component
public class PaymentHandler {
    private final PaymentService paymentService;
    private final ObjectMapper objectMapper;

    public PaymentHandler(PaymentService svc, ObjectMapper om) {
        this.paymentService = svc; this.objectMapper = om;
    }

    public ServerResponse list(ServerRequest request) {
        return ServerResponse.ok().body(paymentService.findAll());
    }

    public ServerResponse getById(ServerRequest request) {
        var id = UUID.fromString(request.pathVariable("id"));
        try {
            return ServerResponse.ok().body(paymentService.findById(id));
        } catch (PaymentService.PaymentNotFoundException e) {
            return ServerResponse.notFound().build();
        }
    }

    public ServerResponse create(ServerRequest request) throws Exception {
        var dto = objectMapper.readValue(
            request.servletRequest().getInputStream(), PaymentDto.class);
        var payment = paymentService.create(dto);
        return ServerResponse
            .created(URI.create("/api/payments/" + payment.getId()))
            .body(payment);
    }
}
\end{lstlisting}

\subsection*{А.26. PaymentRouter.java}
\begin{lstlisting}[caption={PaymentRouter.java --- тип контроллера \textnumero 2: RouterFunction}]
package com.example.idempotency.payment;

import org.springframework.context.annotation.*;
import org.springframework.web.servlet.function.*;
import static org.springframework.http.MediaType.APPLICATION_JSON;
import static org.springframework.web.servlet.function.RequestPredicates.*;

@Configuration
public class PaymentRouter {
    @Bean
    public RouterFunction<ServerResponse> paymentRoutes(PaymentHandler h) {
        return RouterFunctions.route()
            .path("/api/payments", b -> b
                .GET("", accept(APPLICATION_JSON), h::list)
                .GET("/{id}", accept(APPLICATION_JSON), h::getById)
                .POST("", contentType(APPLICATION_JSON), h::create)
            ).build();
    }
}
\end{lstlisting}

% ============================================================
% Домен: Товары (Spring Data REST)
% ============================================================

\subsection*{А.27. Product.java}
\begin{lstlisting}[caption={Product.java --- сущность с @Version для optimistic locking}]
package com.example.idempotency.product;

import jakarta.persistence.*;
import java.math.BigDecimal;

@Entity @Table(name = "products")
public class Product {
    @Id @GeneratedValue(strategy = GenerationType.IDENTITY)
    private Long id;
    @Column(nullable = false, unique = true, length = 100)
    private String sku;
    @Column(nullable = false)
    private String name;
    @Column(nullable = false, precision = 19, scale = 2)
    private BigDecimal price;
    @Column(nullable = false)
    private Integer quantity = 0;
    @Version
    private Integer version;

    public Product() {}
    public Product(String sku, String name, BigDecimal price, Integer qty) {
        this.sku = sku; this.name = name; this.price = price; this.quantity = qty;
    }
    // Getters и Setters
    public Long getId() { return id; }
    public String getSku() { return sku; }
    public String getName() { return name; }
    public void setName(String n) { this.name = n; }
    public Integer getQuantity() { return quantity; }
    public void setQuantity(Integer q) { this.quantity = q; }
    public Integer getVersion() { return version; }
}
\end{lstlisting}

\subsection*{А.28. ProductRepository.java}
\begin{lstlisting}[caption={ProductRepository.java --- тип контроллера \textnumero 3: Spring Data REST}]
package com.example.idempotency.product;

import org.springframework.data.jpa.repository.JpaRepository;
import org.springframework.data.rest.core.annotation.RepositoryRestResource;
import java.util.Optional;

@RepositoryRestResource(collectionResourceRel = "products", path = "products")
public interface ProductRepository extends JpaRepository<Product, Long> {
    Optional<Product> findBySku(String sku);
}
\end{lstlisting}

\subsection*{А.29. ProductEventHandler.java}
\begin{lstlisting}[caption={ProductEventHandler.java --- обработчик событий Spring Data REST}]
package com.example.idempotency.product;

import org.slf4j.*;
import org.springframework.data.rest.core.annotation.*;
import org.springframework.stereotype.Component;

@Component @RepositoryEventHandler
public class ProductEventHandler {
    private static final Logger log =
        LoggerFactory.getLogger(ProductEventHandler.class);

    @HandleBeforeCreate
    public void handleBeforeCreate(Product product) {
        log.info("Creating product: SKU={}", product.getSku());
        if (product.getQuantity() == null) product.setQuantity(0);
    }

    @HandleBeforeSave
    public void handleBeforeSave(Product product) {
        log.info("Updating product: id={}, version={}",
            product.getId(), product.getVersion());
    }
}
\end{lstlisting}

% ============================================================
% Обработка ошибок
% ============================================================

\subsection*{А.30. GlobalExceptionHandler.java}
\begin{lstlisting}[caption={GlobalExceptionHandler.java --- глобальная обработка ошибок}]
package com.example.idempotency.exception;

import com.example.idempotency.order.OrderService;
import com.example.idempotency.payment.PaymentService;
import org.springframework.dao.DataIntegrityViolationException;
import org.springframework.http.*;
import org.springframework.orm.ObjectOptimisticLockingFailureException;
import org.springframework.web.bind.MethodArgumentNotValidException;
import org.springframework.web.bind.annotation.*;
import java.time.LocalDateTime;
import java.util.Map;

@RestControllerAdvice
public class GlobalExceptionHandler {

    @ExceptionHandler(OrderService.OrderNotFoundException.class)
    public ResponseEntity<Map<String, Object>> handleOrderNotFound(
            OrderService.OrderNotFoundException e) {
        return buildError(HttpStatus.NOT_FOUND, e.getMessage());
    }

    @ExceptionHandler(PaymentService.PaymentNotFoundException.class)
    public ResponseEntity<Map<String, Object>> handlePaymentNotFound(
            PaymentService.PaymentNotFoundException e) {
        return buildError(HttpStatus.NOT_FOUND, e.getMessage());
    }

    @ExceptionHandler(DataIntegrityViolationException.class)
    public ResponseEntity<Map<String, Object>> handleDataIntegrity(
            DataIntegrityViolationException e) {
        return buildError(HttpStatus.CONFLICT, "Duplicate resource");
    }

    @ExceptionHandler(ObjectOptimisticLockingFailureException.class)
    public ResponseEntity<Map<String, Object>> handleOptimisticLock(
            ObjectOptimisticLockingFailureException e) {
        return buildError(HttpStatus.CONFLICT, "Version conflict");
    }

    @ExceptionHandler(MethodArgumentNotValidException.class)
    public ResponseEntity<Map<String, Object>> handleValidation(
            MethodArgumentNotValidException e) {
        var msg = e.getBindingResult().getFieldErrors().stream()
            .findFirst().map(err -> err.getField() + ": "
                + err.getDefaultMessage()).orElse("Validation error");
        return buildError(HttpStatus.BAD_REQUEST, msg);
    }

    private ResponseEntity<Map<String, Object>> buildError(
            HttpStatus status, String message) {
        return ResponseEntity.status(status).body(Map.of(
            "timestamp", LocalDateTime.now().toString(),
            "status", status.value(),
            "error", status.getReasonPhrase(),
            "message", message));
    }
}
\end{lstlisting}

% ============================================================
% Тесты
% ============================================================

\subsection*{А.31. BaseIntegrationTest.java}
\begin{lstlisting}[caption={BaseIntegrationTest.java --- базовый класс с Testcontainers}]
package com.example.idempotency;

import org.springframework.boot.test.autoconfigure.web.servlet.AutoConfigureMockMvc;
import org.springframework.boot.test.context.SpringBootTest;
import org.springframework.test.context.*;
import org.testcontainers.containers.PostgreSQLContainer;

@SpringBootTest(webEnvironment = SpringBootTest.WebEnvironment.RANDOM_PORT)
@AutoConfigureMockMvc
@ActiveProfiles("test")
public abstract class BaseIntegrationTest {

    static final PostgreSQLContainer<?> POSTGRES =
        new PostgreSQLContainer<>("postgres:16-alpine")
            .withDatabaseName("idempotency_test")
            .withUsername("test").withPassword("test");

    static { POSTGRES.start(); }

    @DynamicPropertySource
    static void configureProperties(DynamicPropertyRegistry registry) {
        registry.add("spring.datasource.url", POSTGRES::getJdbcUrl);
        registry.add("spring.datasource.username", POSTGRES::getUsername);
        registry.add("spring.datasource.password", POSTGRES::getPassword);
    }
}
\end{lstlisting}

% ============================================================
% Расширение под транспорт-агностичность: V4 + ResponseKind + обновлённые Entity/Service/Filter
% ============================================================

\subsection*{А.32. V4\_\_idempotency\_content\_meta.sql (расширение под произвольный Content-Type)}
\begin{lstlisting}[caption={V4: расширение idempotency\_keys для SOAP/gRPC}, language=SQL]
ALTER TABLE idempotency_keys
    ADD COLUMN response_content_type VARCHAR(150),
    ADD COLUMN response_kind         VARCHAR(16) NOT NULL DEFAULT 'TEXT',
    ADD COLUMN response_body_bytes   BYTEA;
\end{lstlisting}

\subsection*{А.33. ResponseKind.java (дискриминатор тела)}
\begin{lstlisting}[caption={Enum для различения текстового и бинарного снэпшота}]
package com.example.idempotency.idempotency;

public enum ResponseKind { TEXT, BINARY }
\end{lstlisting}

\subsection*{А.34. IdempotencyKeyEntity.java (обновлённая после V4)}
\begin{lstlisting}[caption={JPA-сущность со снэпшотом ответа (extract)}]
@Entity
@Table(name = "idempotency_keys")
public class IdempotencyKeyEntity {
    @Id @GeneratedValue(strategy = GenerationType.IDENTITY)
    private Long id;

    @Column(name = "idempotency_key", nullable = false, unique = true)
    private String idempotencyKey;

    @Column(nullable = false, length = 10)  private String method;
    @Column(nullable = false, length = 500) private String path;

    @Column(name = "response_status")                                   private Integer responseStatus;
    @Column(name = "response_body", columnDefinition = "TEXT")          private String responseBody;
    @Column(name = "response_content_type", length = 150)               private String responseContentType;

    @Enumerated(EnumType.STRING)
    @Column(name = "response_kind", nullable = false, length = 16)
    private ResponseKind responseKind = ResponseKind.TEXT;

    @Basic(fetch = FetchType.LAZY)
    @Column(name = "response_body_bytes", columnDefinition = "BYTEA")
    private byte[] responseBodyBytes;

    @Column(name = "created_at", nullable = false, updatable = false)
    private LocalDateTime createdAt = LocalDateTime.now();
    @Column(name = "expires_at", nullable = false)
    private LocalDateTime expiresAt = LocalDateTime.now().plusHours(24);

    // getters / setters опущены
}
\end{lstlisting}

\subsection*{А.35. IdempotencyService.java (обновлённый с двумя save-методами)}
\begin{lstlisting}[caption={Единая точка интеграции для всех транспортов}]
@Service
public class IdempotencyService {
    private final IdempotencyKeyRepository repository;

    public IdempotencyService(IdempotencyKeyRepository repository) {
        this.repository = repository;
    }

    public Optional<IdempotencyKeyEntity> findByKey(String key) {
        return repository.findByIdempotencyKey(key);
    }

    @Transactional
    public IdempotencyKeyEntity registerKey(String key, String method, String path) {
        return repository.save(new IdempotencyKeyEntity(key, method, path));
    }

    /** Сохранение текстового ответа (REST/SOAP/GraphQL). */
    @Transactional
    public void saveTextResponse(String key, int status, String contentType, String responseBody) {
        repository.findByIdempotencyKey(key).ifPresent(entity -> {
            entity.setResponseStatus(status);
            entity.setResponseContentType(contentType);
            entity.setResponseKind(ResponseKind.TEXT);
            entity.setResponseBody(responseBody);
            repository.save(entity);
        });
    }

    /** Сохранение бинарного ответа (gRPC protobuf). */
    @Transactional
    public void saveBinaryResponse(String key, int status, String contentType, byte[] bytes) {
        repository.findByIdempotencyKey(key).ifPresent(entity -> {
            entity.setResponseStatus(status);
            entity.setResponseContentType(contentType);
            entity.setResponseKind(ResponseKind.BINARY);
            entity.setResponseBodyBytes(bytes);
            repository.save(entity);
        });
    }
}
\end{lstlisting}

\subsection*{А.36. IdempotencyFilter.java (обновлённый, транспорт-агностичный)}
\begin{lstlisting}[caption={Фильтр сохраняет реальный Content-Type и пишет байты через OutputStream}]
public class IdempotencyFilter extends OncePerRequestFilter {
    public static final String IDEMPOTENCY_KEY_HEADER = "Idempotency-Key";
    private static final String JSON_UTF8 = "application/json;charset=UTF-8";

    private final IdempotencyService idempotencyService;
    private final IdempotencyProperties properties;

    @Override
    protected void doFilterInternal(HttpServletRequest request,
                                    HttpServletResponse response,
                                    FilterChain chain) throws ServletException, IOException {
        if (!"POST".equalsIgnoreCase(request.getMethod()) || !properties.isEnabled()) {
            chain.doFilter(request, response); return;
        }
        String key = request.getHeader(IDEMPOTENCY_KEY_HEADER);
        if (key == null || key.isBlank()) {
            writeError(response, HttpStatus.BAD_REQUEST,
                    "{\"error\":\"Idempotency-Key header is required for POST requests\"}");
            return;
        }
        var existing = idempotencyService.findByKey(key);
        if (existing.isPresent()) {
            var cached = existing.get();
            if (cached.getResponseStatus() != null) { replay(response, cached); return; }
            writeError(response, HttpStatus.CONFLICT,
                    "{\"error\":\"Request with this key is already being processed\"}");
            return;
        }
        try { idempotencyService.registerKey(key, request.getMethod(), request.getRequestURI()); }
        catch (DataIntegrityViolationException e) {
            writeError(response, HttpStatus.CONFLICT,
                    "{\"error\":\"Request with this key is already being processed\"}");
            return;
        }
        ContentCachingResponseWrapper wrapper = new ContentCachingResponseWrapper(response);
        chain.doFilter(request, wrapper);
        String contentType = resolveContentType(wrapper);
        String body = new String(wrapper.getContentAsByteArray(), StandardCharsets.UTF_8);
        idempotencyService.saveTextResponse(key, wrapper.getStatus(), contentType, body);
        wrapper.copyBodyToResponse();
    }

    private static void replay(HttpServletResponse response, IdempotencyKeyEntity cached) throws IOException {
        response.setStatus(cached.getResponseStatus());
        String ct = cached.getResponseContentType() != null ? cached.getResponseContentType() : JSON_UTF8;
        response.setContentType(ct);
        byte[] bytes = (cached.getResponseKind() == ResponseKind.BINARY && cached.getResponseBodyBytes() != null)
                ? cached.getResponseBodyBytes()
                : (cached.getResponseBody() != null ? cached.getResponseBody().getBytes(StandardCharsets.UTF_8) : new byte[0]);
        response.setContentLength(bytes.length);
        response.getOutputStream().write(bytes);
        response.getOutputStream().flush();
    }

    private static String resolveContentType(ContentCachingResponseWrapper w) {
        String h = w.getHeader(HttpHeaders.CONTENT_TYPE);
        if (h != null && !h.isBlank()) return h;
        String ct = w.getContentType();
        return (ct != null && !ct.isBlank()) ? ct : JSON_UTF8;
    }
}
\end{lstlisting}

\subsection*{А.37. WebMvcConfig.java (URL-паттерны фильтра для всех четырёх HTTP-транспортов)}
\begin{lstlisting}[caption={Единый бин IdempotencyFilter перехватывает 4 транспорта}]
@Configuration
@EnableConfigurationProperties(IdempotencyProperties.class)
public class WebMvcConfig {
    @Bean
    public FilterRegistrationBean<IdempotencyFilter> idempotencyFilterRegistration(
            IdempotencyService svc, IdempotencyProperties props) {
        var reg = new FilterRegistrationBean<IdempotencyFilter>();
        reg.setFilter(new IdempotencyFilter(svc, props));
        reg.addUrlPatterns(
                "/api/orders/*", "/api/orders",
                "/api/payments/*", "/api/payments",
                "/api/soap/*",
                "/graphql");
        reg.setOrder(1);
        return reg;
    }
}
\end{lstlisting}

% ============================================================
% SOAP-транспорт
% ============================================================

\subsection*{А.38. V5\_\_shipments.sql + xsd/shipments.xsd (SOAP-домен)}
\begin{lstlisting}[caption={V5 миграция и XSD контракт SOAP}, language=SQL]
-- V5__shipments.sql
CREATE TABLE shipments (
    id              UUID PRIMARY KEY DEFAULT gen_random_uuid(),
    tracking_number VARCHAR(64) NOT NULL,
    recipient       VARCHAR(255) NOT NULL,
    status          VARCHAR(32) NOT NULL DEFAULT 'CREATED',
    created_at      TIMESTAMP NOT NULL DEFAULT NOW(),
    CONSTRAINT uk_shipment_tracking_number UNIQUE (tracking_number)
);
\end{lstlisting}
\begin{lstlisting}[caption={xsd/shipments.xsd}, language=XML]
<?xml version="1.0" encoding="UTF-8"?>
<xs:schema xmlns:xs="http://www.w3.org/2001/XMLSchema"
           xmlns:tns="http://example.com/idempotency/shipments"
           targetNamespace="http://example.com/idempotency/shipments"
           elementFormDefault="qualified">
    <xs:element name="createShipmentRequest">
        <xs:complexType><xs:sequence>
            <xs:element name="trackingNumber" type="xs:string"/>
            <xs:element name="recipient" type="xs:string"/>
        </xs:sequence></xs:complexType>
    </xs:element>
    <xs:element name="createShipmentResponse">
        <xs:complexType><xs:sequence>
            <xs:element name="id" type="xs:string"/>
            <xs:element name="trackingNumber" type="xs:string"/>
            <xs:element name="status" type="xs:string"/>
        </xs:sequence></xs:complexType>
    </xs:element>
</xs:schema>
\end{lstlisting}

\subsection*{А.39. ShipmentEndpoint.java + WebServiceConfig.java}
\begin{lstlisting}[caption={SOAP-эндпоинт и регистрация MessageDispatcherServlet}]
@Endpoint
public class ShipmentEndpoint {
    public static final String NAMESPACE = "http://example.com/idempotency/shipments";
    private final ShipmentService service;

    @PayloadRoot(namespace = NAMESPACE, localPart = "createShipmentRequest")
    @ResponsePayload
    public DOMSource createShipment(@RequestPayload DOMSource request) throws Exception {
        Element root = rootElement(request);
        Shipment created = service.create(childText(root, "trackingNumber"),
                                           childText(root, "recipient"));
        Document doc = newDocument();
        Element resp = doc.createElementNS(NAMESPACE, "createShipmentResponse");
        doc.appendChild(resp);
        appendChild(doc, resp, "id", created.getId().toString());
        appendChild(doc, resp, "trackingNumber", created.getTrackingNumber());
        appendChild(doc, resp, "status", created.getStatus());
        return new DOMSource(doc);
    }
    // rootElement / childText / appendChild / newDocument опущены
}

@EnableWs @Configuration
public class WebServiceConfig extends WsConfigurerAdapter {
    @Bean
    public ServletRegistrationBean<MessageDispatcherServlet> messageDispatcherServlet(
            ApplicationContext ctx) {
        MessageDispatcherServlet servlet = new MessageDispatcherServlet();
        servlet.setApplicationContext(ctx);
        servlet.setTransformWsdlLocations(true);
        return new ServletRegistrationBean<>(servlet, "/api/soap/*");
    }
    @Bean(name = "shipments")
    public DefaultWsdl11Definition shipmentsWsdl(XsdSchema shipmentsSchema) {
        DefaultWsdl11Definition wsdl = new DefaultWsdl11Definition();
        wsdl.setPortTypeName("ShipmentsPort");
        wsdl.setLocationUri("/api/soap");
        wsdl.setTargetNamespace(ShipmentEndpoint.NAMESPACE);
        wsdl.setSchema(shipmentsSchema);
        return wsdl;
    }
    @Bean
    public XsdSchema shipmentsSchema() {
        return new SimpleXsdSchema(new ClassPathResource("xsd/shipments.xsd"));
    }
}
\end{lstlisting}

% ============================================================
% gRPC-транспорт
% ============================================================

\subsection*{А.40. notification.proto + V6\_\_notifications.sql}
\begin{lstlisting}[caption={Контракт gRPC и таблица домена}, language={}]
// src/main/proto/notification.proto
syntax = "proto3";
package com.example.idempotency.grpc.notification;
option java_multiple_files = true;
option java_package = "com.example.idempotency.grpc.notification.proto";

service NotificationService {
  rpc Send(SendRequest) returns (SendResponse);
  rpc List(ListRequest) returns (ListResponse);
}
message SendRequest  { string recipient = 1; string message = 2; }
message SendResponse { string id = 1; string status = 2; }
message ListRequest  {}
message ListResponse { repeated NotificationItem items = 1; }
message NotificationItem {
  string id = 1; string recipient = 2; string message = 3; string status = 4;
}
\end{lstlisting}
\begin{lstlisting}[caption={V6\_\_notifications.sql}, language=SQL]
CREATE TABLE notifications (
    id            UUID PRIMARY KEY DEFAULT gen_random_uuid(),
    recipient     VARCHAR(255) NOT NULL,
    message       TEXT NOT NULL,
    message_hash  VARCHAR(64) NOT NULL,
    status        VARCHAR(32) NOT NULL DEFAULT 'SENT',
    created_at    TIMESTAMP NOT NULL DEFAULT NOW(),
    CONSTRAINT uk_notification_recipient_hash UNIQUE (recipient, message_hash)
);
\end{lstlisting}

\subsection*{А.41. GrpcIdempotencyInterceptor.java (главная часть работы)}
\begin{lstlisting}[caption={Transport-generic перехватчик: переиспользует IdempotencyService}]
@GrpcGlobalServerInterceptor
public class GrpcIdempotencyInterceptor implements ServerInterceptor {
    private static final Metadata.Key<String> IDEMPOTENCY_KEY =
            Metadata.Key.of("idempotency-key", Metadata.ASCII_STRING_MARSHALLER);
    private static final String CONTENT_TYPE = "application/grpc+proto";

    private final IdempotencyService idempotencyService;
    private final IdempotencyProperties properties;

    @Override
    public <ReqT, RespT> ServerCall.Listener<ReqT> interceptCall(
            ServerCall<ReqT, RespT> call, Metadata headers,
            ServerCallHandler<ReqT, RespT> next) {

        if (!properties.isEnabled()) return next.startCall(call, headers);

        String key = headers.get(IDEMPOTENCY_KEY);
        if (key == null || key.isBlank()) {
            call.close(Status.INVALID_ARGUMENT.withDescription(
                    "idempotency-key metadata is required"), new Metadata());
            return noopListener();
        }
        var existing = idempotencyService.findByKey(key);
        if (existing.isPresent()) {
            IdempotencyKeyEntity cached = existing.get();
            if (cached.getResponseStatus() != null && cached.getResponseBodyBytes() != null) {
                try {
                    RespT msg = call.getMethodDescriptor().getResponseMarshaller()
                            .parse(new ByteArrayInputStream(cached.getResponseBodyBytes()));
                    call.sendHeaders(new Metadata());
                    call.sendMessage(msg);
                    call.close(Status.OK, new Metadata());
                } catch (Exception e) {
                    call.close(Status.INTERNAL.withDescription("replay failed: " + e.getMessage()),
                            new Metadata());
                }
                return noopListener();
            }
            call.close(Status.ABORTED, new Metadata());
            return noopListener();
        }
        try { idempotencyService.registerKey(key, "GRPC", call.getMethodDescriptor().getFullMethodName()); }
        catch (DataIntegrityViolationException e) {
            call.close(Status.ALREADY_EXISTS, new Metadata());
            return noopListener();
        }
        return next.startCall(new CachingServerCall<>(call, key), headers);
    }

    private static <ReqT> ServerCall.Listener<ReqT> noopListener() {
        return new ServerCall.Listener<>() {};
    }

    private final class CachingServerCall<ReqT, RespT>
            extends ForwardingServerCall.SimpleForwardingServerCall<ReqT, RespT> {
        private final String key;
        CachingServerCall(ServerCall<ReqT, RespT> delegate, String key) {
            super(delegate); this.key = key;
        }
        @Override
        public void sendMessage(RespT message) {
            try {
                byte[] bytes = getMethodDescriptor().getResponseMarshaller()
                        .stream(message).readAllBytes();
                idempotencyService.saveBinaryResponse(key,
                        Status.OK.getCode().value(), CONTENT_TYPE, bytes);
            } catch (IOException e) { /* log only */ }
            super.sendMessage(message);
        }
    }
}
\end{lstlisting}

\subsection*{А.42. NotificationGrpcService.java}
\begin{lstlisting}[caption={Реализация gRPC-сервиса (@GrpcService)}]
@GrpcService
public class NotificationGrpcService
        extends NotificationServiceGrpc.NotificationServiceImplBase {
    private final NotificationService service;
    private final NotificationRepository repository;

    @Override
    public void send(SendRequest req, StreamObserver<SendResponse> observer) {
        Notification n = service.send(req.getRecipient(), req.getMessage());
        observer.onNext(SendResponse.newBuilder()
                .setId(n.getId().toString()).setStatus(n.getStatus()).build());
        observer.onCompleted();
    }

    @Override
    public void list(ListRequest req, StreamObserver<ListResponse> observer) {
        ListResponse.Builder b = ListResponse.newBuilder();
        repository.findAll().forEach(n -> b.addItems(NotificationItem.newBuilder()
                .setId(n.getId().toString()).setRecipient(n.getRecipient())
                .setMessage(n.getMessage()).setStatus(n.getStatus()).build()));
        observer.onNext(b.build());
        observer.onCompleted();
    }
}
\end{lstlisting}

% ============================================================
% GraphQL-транспорт
% ============================================================

\subsection*{А.43. schema.graphqls + V7\_\_inventories.sql}
\begin{lstlisting}[caption={GraphQL-схема (SDL) и таблица домена}, language={}]
# src/main/resources/graphql/schema.graphqls
type Inventory { id: ID!, itemName: String!, quantity: Int! }
type Query { inventories: [Inventory!]!, inventory(id: ID!): Inventory }
type Mutation { reserveInventory(input: ReserveInput!): Inventory! }
input ReserveInput { itemName: String!, quantity: Int! }
\end{lstlisting}
\begin{lstlisting}[caption={V7\_\_inventories.sql}, language=SQL]
CREATE TABLE inventories (
    id          UUID PRIMARY KEY DEFAULT gen_random_uuid(),
    item_name   VARCHAR(255) NOT NULL,
    quantity    INTEGER NOT NULL,
    created_at  TIMESTAMP NOT NULL DEFAULT NOW(),
    CONSTRAINT uk_inventory_item_name UNIQUE (item_name)
);
\end{lstlisting}

\subsection*{А.44. InventoryGraphqlController.java}
\begin{lstlisting}[caption={GraphQL-контроллер с @QueryMapping / @MutationMapping}]
@Controller
public class InventoryGraphqlController {
    private final InventoryRepository repository;

    @QueryMapping
    public List<Inventory> inventories() { return repository.findAll(); }

    @QueryMapping
    public Inventory inventory(@Argument String id) {
        return repository.findById(UUID.fromString(id)).orElse(null);
    }

    @MutationMapping
    @Transactional
    public Inventory reserveInventory(@Argument ReserveInput input) {
        return repository.save(new Inventory(input.itemName(), input.quantity()));
    }

    public record ReserveInput(String itemName, int quantity) {}
}
\end{lstlisting}

\subsection*{А.45. V8\_\_drop\_new\_constraints.sql (профиль non-idempotent)}
\begin{lstlisting}[caption={Дроп UNIQUE на трёх новых доменах --- демонстрация багов}, language=SQL]
ALTER TABLE shipments     DROP CONSTRAINT IF EXISTS uk_shipment_tracking_number;
ALTER TABLE notifications DROP CONSTRAINT IF EXISTS uk_notification_recipient_hash;
ALTER TABLE inventories   DROP CONSTRAINT IF EXISTS uk_inventory_item_name;
\end{lstlisting}
